\documentclass[10pt]{article}
\usepackage[margin=20mm]{geometry}
\usepackage{amsmath,amssymb,mathtools,fontspec,unicode-math,xurl,hyperref}
\IfFontExistsTF{DejaVu Serif}{\setmainfont{DejaVu Serif}}{\setmainfont{Latin Modern Roman}}
\setmathfont{latinmodern-math.otf}
\setlength{\parindent}{0pt}\setlength{\parskip}{6pt}\sloppy
\title{Technical review of Rigidity of steady solutions to the Navier–Stokes equations in high dimensions and its applications}
\author{AI-assisted technical review}
\date{}
\begin{document}
\maketitle
Abstract: The bound arXiv v2 manuscript is reviewed through the weighted rigidity proof and its singularity and exterior-decay applications. An admissible exact radial solution disproves the forcing identity in the nonzero-flux step of the exterior-decay proof. This is assessed as Major for that application proof; adding the missing interaction term repairs the identity without worsening the stated estimates. A separate half-space flux-sign inconsistency is Minor. No counterexample to the main rigidity or decay conclusions is asserted. Actual image inspection remains blocked.\par
\section*{Binding and operative hypotheses}
The authors are Jeaheang Bang, Changfeng Gui, Hao Liu, Yun Wang, and Chunjing Xie. The exact reviewed source is arXiv:2306.05184v2, dated 2 September 2023 [paper]. It is a version-limited preprint associated with DOI 10.4171/jems/1738 (https://doi.org/10.4171/jems/1738). The publisher’s abstract and dates differ from this preprint record [publisher]; no transfer of the findings to an unexamined publisher file is made.\par
All 26 extracted text pages were actually read. The equations are \[-\Delta u+(u\cdot\nabla)u+\nabla p=0,\qquad \operatorname{div}u=0,\] with smoothness off the origin and the global scale bound \(|u(x)|\leq C/|x|\). The full-space theorem and Navier-slip half-space theorem use \(n\geq4\); the Dirichlet half-space assertion uses \(n=4\) only. Neither smallness nor any self-similarity is assumed for Theorem 1.1. Smallness is introduced later in Lemma 4.2, after a blowdown argument improves exterior decay to \(o(|x|^{-1})\).\par
The scientific review distinguishes the full-space rigidity result, removal of the isolated singularity, and the general-flux exterior application. A defect in the last proof is not evidence that the earlier theorem is false.\par
\section*{P1: the nonzero-flux forcing is incorrect}
\subsection*{Precise identity and an admissible solution}
On page 23, Step 4 of Lemma 4.2 sets \[a(x)=c\frac{x}{|x|^n},\qquad p_a=-\frac12|a|^2,\qquad u_1=u-a.\] Equation (73) is asserted for the divergence-free extension \(\widetilde u_1\), with pressure \(\eta p-\eta p_a\). The displayed forcing ends with \[-(\eta u\cdot\nabla)a-(\eta a\cdot\nabla)u.\] These are not the correct interaction terms for \(u_1\). Before the cutoff is introduced, subtraction of the two Navier–Stokes equations gives \[\begin{aligned}
-\Delta u_1+(u_1\cdot\nabla)u_1+\nabla(p-p_a)
 & =-(u_1\cdot\nabla)a-(a\cdot\nabla)u_1\\
 & =-(u\cdot\nabla)a-(a\cdot\nabla)u
       +2(a\cdot\nabla)a.\end{aligned}\] The last term is missing from the printed forcing.\par
There is an exact countercheck satisfying the exterior lemma’s hypotheses. Take any \(n\geq4\), \(R=1\), and \[u=a,\qquad p=p_a,\qquad 0<|c|\leq\epsilon.\] This is smooth in the exterior domain, has nonzero flux \(c|S^{n-1}|\), and satisfies \[\operatorname{div}a=0,\quad \Delta a=0,\quad
(a\cdot\nabla)a=(1-n)c^2\frac{x}{|x|^{2n}},\quad
\nabla p_a=-(a\cdot\nabla)a.\] Consequently it solves the exact steady equations. Also \[|a|=|c|\,|x|^{1-n}\leq\epsilon/|x|\quad (|x|>1).\] The parameter can be arbitrarily small, so no unknown smallness threshold is bypassed.\par
For this solution \(u_1=0\). Outside \(B_3\), \(\eta=1\) and the compactly supported divergence correction vanishes. The left-hand side of (73) is zero, whereas its printed right-hand side is \[-2(a\cdot\nabla)a=2(n-1)c^2\frac{x}{|x|^{2n}}\ne0.\] In dimension four, at \(x=(4,0,0,0)\) its first component is \(3c^2/8192\). The reviewer script independently verifies divergence, Laplacian, pressure cancellation, and this residual. Its numerical illustration at \(c=1/100\) gives \(3/81920000\); the on-hypothesis countercheck is the parametric one \(0<|c|\leq\epsilon\), not an assumption that \(\epsilon\geq1/100\).\par
\subsection*{Impact and verified correction}
The assessment is Major for the application proof. The extension to arbitrary nonzero flux, needed for Theorem 1.3 as stated, uses an incorrectly identified whole-space equation. The zero-flux proof and the preceding rigidity theorem do not depend on this formula. The radial solution also decays faster than the claimed optimal upper bound, so it is not a counterexample to Theorem 1.3.\par
Add \[2\eta(a\cdot\nabla)a\] to the displayed forcing, or write the two interactions directly using \(u_1\). This is the exact algebraic correction. On the exterior, \[|2(a\cdot\nabla)a|=2(n-1)c^2|x|^{1-2n}
=O\bigl(|x|^{-n-1}\bigr)\] for \(n\geq4\), with strictly faster decay than needed. The cutoff removes the singularity from its support, and the coefficient is quadratic in the small flux parameter. Thus the corrected forcing satisfies the same decay and smallness bounds used in the subsequent Stokes convolution and contraction argument. This local repair is verified; no author adoption or complete independent reproving of the cited analytic machinery is claimed.\par
\section*{P2: half-space energy flux has the opposite sign}
Equation (51) on page 17 states that the half-annular Dirichlet energy equals \(h(R_2)-h(R_1)\), but defines \(h\) with a minus sign in front of the integral of \(\partial_r u\cdot u-\tfrac12|u|^2u^r-pu^r\). Let \[F(R)=\int_{(\partial B_R)_+}
\left(\partial_r u\cdot u-Hu^r\right)\,d\sigma,
\qquad H=\frac12|u|^2+p.\] Direct multiplication of the momentum equation by \(u\) and integration by parts gives \[\int_{B_{R_2,+}\setminus B_{R_1,+}}|\nabla u|^2
=F(R_2)-F(R_1).\] The flat boundary terms vanish by the Dirichlet condition. With the printed definition \(h=-F\), the right-hand side is therefore \(h(R_1)-h(R_2)\), and \(h\) is decreasing.\par
An independent local check uses the exact Couette solution \(u=(b x_4,0,0,0)\), \(p=0\). On a four-dimensional half-annulus, its energy is \(b^2\pi^2(R_2^4-R_1^4)/4\) and \(F(R)=b^2\pi^2R^4/4\); the convective boundary contribution integrates to zero by oddness in \(x_1\). This checks the local integration formula. This Couette field does not satisfy the theorem’s global \(C/|x|\) bound and is not presented as a counterexample to the rigidity conclusion.\par
The assessment is Minor. Removing the minus in the definition of \(h\), or reversing both the difference and monotonicity, restores consistency. Boundedness, finite energy, Sobolev integrability, the vanishing endpoint subsequences, and the final zero-energy conclusion all survive. The calculation concerns the derivation and sign convention, not a claimed nonzero globally decaying half-space solution.\par
\section*{Checks of the main argument}
\subsection*{Pressure normalization and four dimensions}
Lemma 2.1 first obtains derivative bounds on annuli. The averages of \(p\) over dyadic spheres differ by \(O(2^{-2k})\), so they converge at infinity. Subtracting that single limiting constant, followed by the gradient estimate along paths within comparable-radius annuli, gives \(|p(x)|\leq C|x|^{-2}\) globally. This normalization is necessary before assigning a sign to \(H\); pressure is not arbitrarily renormalized again during the proof.\par
For \(n=4\), equation (27) loses its two \((n-4)\) terms: \[\int_{B_{R_2}\setminus B_{R_1}}|\nabla u|^2
=h_1(R_2)-h_1(R_1).\] The boundary estimates make \(h_1\) bounded, and it is increasing. Hence the full Dirichlet integral is finite. The puncture does not create a first-derivative distributional defect: the boundary scale of \(u\) is \(O(\varepsilon^2)\) in dimension four. Together with decay at infinity, the homogeneous Sobolev inequality gives \(u\in L^4\).\par
Lemma A.1 supplies sequences at both radial endpoints for which \(R\int_{\partial B_R}|u|^4\to0\). It supplies a lower-limit statement, not full pointwise convergence. The estimate \(|h_1(R)|\leq C(R\int_{\partial B_R}|u|^4)^{1/4}\) and monotonicity then force both actual limits of \(h_1\) to zero. Thus the energy vanishes. This uses no smallness assumption.\par
\subsection*{The positive head pressure and higher dimensions}
For \(n\geq5\), equation (33) is tested against \(H_+^\alpha\), where \(\alpha=(n-4)/2>0\). The boundary functional \(h_2\) is decreasing and scale bounded. The resulting finite integral of \(|\nabla H_+|^2H_+^{\alpha-1}\) gives \(\nabla(H_+^{(\alpha+1)/2})\in L^2\). The Sobolev exponent is exactly \[\frac{\alpha+1}{2}\frac{2n}{n-2}=\frac n2.\] The subsequent boundary estimate is a power \(1-4/n>0\) of \(R\int_{\partial B_R}H_+^{n/2}\). Endpoint subsequences and monotonicity therefore yield \(H_+=0\). The fractional-power test at zeros can be justified by positive regularization; it is not a maximum-principle argument assuming boundedness at the singular point.\par
The next multiplier is \(\nabla r^{4-n}\). Its Laplacian-velocity flux vanishes because \(\operatorname{div}\Delta u=0\) and the small-sphere estimate is \(O(\varepsilon^{n-4})\to0\). This is where \(n\geq5\) is essential. The resulting increasing, bounded \(h_3\) controls \[\int_{\mathbb R^n}\bigl((n-2)|u^r|^2-2H\bigr)r^{2-n}<\infty.\] Since \(H\leq0\), the integrands have the required signs. This yields weighted integrability of \(H\), then of \(|u|^2\) and \(|\nabla u|^2\) through (27). The pressure estimate needed here follows directly from \(|p|\leq|H|+|u|^2/2\).\par
The integrability of \(|h_3(R)|/R\) gives endpoint subsequences tending to zero. Monotonicity yields \(H=u^r=0\); the remaining weighted energy identity then forces \(u=0\). Both radial endpoints and the dependence on \(H\leq0\) are retained in this reconstruction.\par
\subsection*{Boundary conditions and applications}
The Navier-slip condition on the flat boundary permits even reflection of tangential velocity and pressure and odd reflection of the normal velocity. The normal-velocity trace and tangential normal-derivative conditions prevent interface defects. Local regularity then reduces this case to the full-space theorem. The no-slip case uses only the four-dimensional energy method, with the sign correction described above. A no-slip theorem in all higher dimensions is explicitly not claimed.\par
The singularity-removal argument uses a cited smallness criterion to produce a nonzero value on a rescaled unit sphere if regularity fails. Uniform local estimates on compact punctured sets yield a global limiting solution, to which Theorem 1.1 applies; the nonzero normalized value contradicts that limit. The exterior blowdown uses the same compactness mechanism to obtain \(o(|x|^{-1})\).\par
For the final decay improvement, the zero-flux cutoff admits a Bogovskii correction. The Stokes kernel bounds lead to a contraction in a weighted supremum class with decay \(|x|^{2-n}\); uniqueness is estimated in the weaker \(|x|^{-1}\) class. The definition in the manuscript imposes \(C^1\) regularity while using only a supremum norm. A precise fixed-point formulation should use the complete weighted continuous, distributionally divergence-free space first and recover smoothness by the cited local Stokes regularity. No automatic completeness of the literal \(C^1\) space is assumed here. The nonzero-flux step requires P1’s corrected forcing.\par
\section*{Outcome and limitations}
The principal supported finding is the incorrect nonzero-flux equation, assessed as Major for its role in a main application proof. Its corrected interaction retains the required estimates. The half-space flux sign is a separate Minor inconsistency. No nontrivial global \(C/|x|\) solution in the theorem’s dimension range has been produced.\par
External Stokes estimates, Bogovskii construction, and the cited removable-singularity criterion remain dependencies. The publisher and arXiv records, including bounded searches for associated corrections and supplementary code, did not identify a separate package for this exact preprint. The arXiv source link is a manuscript source, not evidence of an executable experiment. No author code was run.\par
Item | Assessment | Supported scope\par
P1 | Major | Nonzero-flux forcing; local repair verified\par
P2 | Minor | Half-space flux sign; endpoint argument survives\par
Core rigidity | No defect established | Weighted proof inspected with cited analytic inputs\par
Scientific decision | Pending | No independent or human acceptance\par
Delivery | Incomplete | Actual source/report image inspection missing\par
\section*{Supporting checks}
Paper ID: `M25-0359`.\par
> 00002TOKEN.\par
Reasoning and writing are by this assistant session. Native tools perform mechanical calculations, compilation, parsing, and rendering; they are not independent reviewers. Controller UI observations are attributed to the controller, and backend equivalence is unverified. The actual image-reading interface returns `EMFILE`. No rendering operation is counted as viewing.\par
9 J. Bang, C. Gui, H. Liu, Y. Wang and C. Xie, Rigidity of Steady Solutions to the Navier–Stokes Equations in High Dimensions and its applications, arXiv:2306.05184v2, 2 September 2023, https://arxiv.org/abs/2306.05184v2. EMS Press, DOI 10.4171/jems/1738, https://ems.press/journals/jems/articles/14299214, consulted 27 October 2026.\par
\end{document}
